\documentclass[11pt,a4paper]{article}
\usepackage[T1]{fontenc}
\usepackage{lmodern,amsmath,amssymb,amsthm,booktabs,microtype}
\usepackage[margin=27mm]{geometry}
\usepackage[colorlinks=true,linkcolor=blue,urlcolor=blue]{hyperref}
\setlength{\emergencystretch}{2em}
\newtheorem{theorem}{Theorem}[section]
\newtheorem{proposition}[theorem]{Proposition}
\newtheorem{lemma}[theorem]{Lemma}
\DeclareMathOperator{\rad}{rad}
\newcommand{\Q}{\mathbb Q}\newcommand{\Z}{\mathbb Z}
\title{Single-cover criteria from the signed mixed-power height bound}
\author{}\date{17 November 2026}
\begin{document}\maketitle

Assume the weak rational $abc$ inequality
\begin{equation}\label{eq:abc}
 c\le K_\kappa\rad(abc)^\kappa
\end{equation}
for positive coprime integers $a+b=c$, with $\kappa>1$ and $K_\kappa\ge1$.

\section{Source and height constants}
Retain the conditional single-cover source of \cite{Single}, its
three point roles $A,H,U$, its guarded finite point algebra, and all
of its group, formal-comparison and trace-height premises. Write
$A=[n]P$, where $n\in\Z\setminus\{0\}$, and set
\[
 \beta=X^2+B,\qquad T=\operatorname{Tr}_{K/\Q}x(A),\qquad
 H_0=\beta^e+T^e,\qquad t=H_0^d+z^2,
\]
where $B\ge3$ is fixed and odd, $e$ is a positive integer, $d$ is
a positive even integer,
and $z\in\Q$ is the sole free detector parameter.
There are no new coordinates for $H_0$ or $t$.

Assume the fixed weak rational $abc$ inequality with exponent
$\kappa>1$ and constant $K_\kappa\ge1$, and put
\[
 \epsilon=d-\kappa(d/2+2),\quad
 \alpha_e=e/\kappa-3,\quad
 \gamma=8/\kappa-6.
\]
All three quantities are required to be positive. The direct point
chart of \cite[Section 3]{Conditional} has
$x(A)=(h_0+h_1\sqrt2)^2$, so
\begin{equation}\label{eq:positiveT}
 T=2(h_0^2+2h_1^2)>0.
\end{equation}
The nonzero point guard supplies strictness. This rational trace is
independent of every finite ordinate-root choice. Hence $\beta,T$ are
positive at every point-algebra character, and $e$ need not be even.

\begin{lemma}\label{lem:positivepowers}
Let $a,b>0$ be rational. For any positive integer $e$ with
$\alpha=e/\kappa-3>0$, one has
\[
 h(1:a:b)\le\alpha^{-1}h(a^e+b^e)
 +\frac{\alpha+1}{e\alpha}\log2
 +\frac{\log K_\kappa}{\kappa\alpha}.
\]
\end{lemma}
\begin{proof}
Use the denominator normalization of \cite[Lemma 2.2]{Conditional}:
write $a=gA/D$, $b=gB/D$, with positive integers $A,B,g,D$,
$\gcd(A,B)=\gcd(g,D)=1$. Put $H=\max(A,B)$ and
$N=A^e+B^e$, so $H^e\le N\le2H^e$. Write
\[
 k_0=\gcd(N,D^e),\quad m_0=N/k_0,\quad n_0=D^e/k_0,
 \quad m=g^em_0,\quad M=\max(m,n_0).
\]
The fraction $m/n_0$ is reduced. Every prime in $N$ divides $D$
or survives in $m_0$, whence $\rad(N)\le Dm_0$. The positive
coprime triple $A^e+B^e=N$ and \eqref{eq:abc} imply
\[
 H^{e/\kappa-2}\le K_\kappa^{1/\kappa}Dm_0,
 \qquad
 H^\alpha\le2^{1/e}K_\kappa^{1/\kappa}
               m_0^{(e-1)/e}n_0^{1/e}.
\]
Also $\max(gH,D)\le2^{1/e}H(M/m_0)^{1/e}$.
Combining logarithms bounds $h(1:a:b)$ by the stated additive
constant plus
\[
 \frac{e-1-\alpha}{e\alpha}\log m_0
 +\frac1{e\alpha}\log n_0+\frac1e\log M.
\]
The three coefficients are nonnegative and sum to $1/\alpha$;
since $m_0,n_0\le M$, the assertion follows. Positivity of $a,b$
supplies every sign assertion in this argument for both parities of $e$.
\end{proof}

Let $\mathcal M(t)$ be the retained odd contact radical, and write
$E_h=h(1:\beta:T)$. The source's contact estimate and
Lemma~\ref{lem:positivepowers} give
\[
 h(t)\le\gamma^{-1}\log\mathcal M(t)+O(1),\qquad
 E_h\le\alpha_e^{-1}h(H_0)+O(1).
\]
\cite[Theorem 1.1]{Height}, applied to the plus sign, gives
\begin{equation}\label{eq:height}
 E_h\le q\log\mathcal M(t)+O(1),\qquad
 q=\frac{2\kappa-1}{\epsilon\alpha_e\gamma}.
\end{equation}
Every constant is independent of $B$, $X$ and the source tuple.

If the argument is written $H_0^{dm}+z^2$, apply
\cite[Theorem 1.1]{Height} directly
with the single even exponent $D=dm$. For $m>1$ this is stronger
than first applying the degree-$d$ bound to $H_0^m$, because
\[
 \epsilon_{dm}=m\epsilon_d+2\kappa(m-1)>m\epsilon_d.
\]
The optimization below therefore uses one total even exponent $d$.

\section{Fixed and varying square classes}

Retain the native single quadratic radicand $\mathcal R$ of
\cite{Single}. For $k=0$ allow its rational-square condition, or
the established fixed finite positive $S$-unit class menu. For
positive even $k$, retain the varying-class construction of
\cite{Classes}: its positive class has odd support satisfying
\[
 \log\operatorname{rad}_{S}^{\mathrm{odd}}(\delta)
       \le k E_h+O(1).
\]
In particular the finite menu
\[
 \{1,\beta^2+T^2,1+X^4\}
\]
has $k=2$. The menu can also be multiplied by any fixed finite
positive $S$-unit menu. All its previous preparation and
whole-character guards are retained.

\begin{theorem}\label{thm:sound}
The single-cover source is sound for every sufficiently large fixed
odd $B$, uniformly in all its rational witnesses, whenever
\begin{equation}\label{eq:absorption}
 q(k+1/2)<1.
\end{equation}
For $k=0$ this is the condition $q<2$.
\end{theorem}
\begin{proof}
Suppose $X\ne n$. The source's local extraction and the support
budget of \cite{Classes} give
\[
 \log\mathcal M(t)
 \le(k+1/2)E_h+\log(1+|n|)+O(1).
\]
Here the factor $1/2$ uses the exact input inequality
$2h(X)\le h(\beta)$, valid for $\beta=X^2+B$ with $B\ge1$.
Substitution in \eqref{eq:height} and strict absorption give
\[
 E_h\le C\log(1+|n|)+C'
\]
with constants independent of $B$. The retained trace-height theorem
then bounds $n^2$ by the same type of expression, so both $n$ and
$E_h$ are bounded independently of $B$. This contradicts
$E_h\ge h(\beta)\ge\log B$ for all sufficiently large fixed
odd $B$. Thus $X=n\in\Z$.
\end{proof}

\section{Exact parameter rows and argument degrees}

For the fixed-class case $k=0$, the following rows have strict
contraction:
\[
\begin{array}{c|c|c|c|c|c}
\kappa&e&d&2ed&q&q/2\\\hline
101/100&4&6&48&520251/893855&520251/1787710\\
25/24&4&6&48&8125/8379&8125/16758\\
16/15&4&6&48&68/45&34/45\\
121/100&6&8&96&1039511/648906&1039511/1297812
\end{array}
\]
For the varying menu with $k=2$, exact rows are
\[
\begin{array}{c|c|c|c|c|c}
\kappa&e&d&2ed&q&(5/2)q\\\hline
101/100&5&6&60&520251/1815355&520251/726142\\
25/24&6&6&72&8125/27531&40625/55062\\
16/15&5&8&80&68/243&170/243\\
121/100&7&12&168&1039511/2892808&5197555/5785616
\end{array}
\]
Ordinary rational $abc$ supplies the first row of each table.
The remaining rows use only the indicated fixed weak exponents,
together with the retained source premises.

The column $2ed$ is the exact degree of the scalar argument $t$ in
the retained input, point and square-parameter coordinates. Indeed
$\beta=X^2+B$ has degree two, and in the direct square-abscissa
chart $T=2(h_0^2+2h_1^2)$ also has degree two. Hence $H_0$ has
degree $2e$ and $t=H_0^d+z^2$ has degree $2ed$. This column is
not a degree bound for the full defining determinant.

\begin{proposition}\label{prop:optimal}
Within the displayed power criterion, the least argument degree
under ordinary rational $abc$ is $48$ for $k=0$ and $60$ for $k=2$.
For each fixed exponent in the two tables, its row has the smallest
argument degree among all integers $e$ and even integers $d$ satisfying
the stated positivity and contraction inequalities. Among pairs of
the same minimal degree the table chooses the smallest $d$.
\end{proposition}
\begin{proof}
Positivity with $\kappa>1$ forces $e\ge4$ and $d\ge6$.
Thus $2ed\ge48$, attained in the first table.
For $k=2$, note
\[
 \frac{\partial}{\partial\kappa}
 \frac{2\kappa-1}{d-\kappa(d/2+2)}
 =\frac{3d/2-2}{\epsilon^2}>0.
\]
Both $\alpha_e$ and $\gamma$ decrease with $\kappa$, so $q$
increases on its positive domain. Its best limiting value at
$\kappa=1$ is
\[
 q_{1,e,d}=\frac1{(d-4)(e-3)}.
\]
The only pair with integer $e\ge4$, even $d\ge6$, and $2ed<60$ is
$(e,d)=(4,6)$. At this pair the limiting contraction is $5/4>1$.
Degree $60$ is attained by the first row of the second table.

For each fixed exponent, a smaller argument degree would require
$2ed$ smaller than the degree of its displayed row. There are only
finitely many such pairs, since $e\ge4,d\ge6$. The accompanying
exact rational certificate lists every one, checks the positivity
conditions, and records a contraction at least one for each remaining
pair. It also checks all rows and all pairs of the same degree
directly. Thus no unbounded search
range or numerical approximation is used for this assertion.
\end{proof}

\section{Guards, coordinates and the remaining realization problem}

No detector, pole-extraction, finite-algebra, or guard equation is
changed in this height argument. At every prepared real character,
$H_0>1$ and $t>1$, so the previous positivity and radicand-unit
proofs still apply for every row. The scalar shift $B$ enters only
through $\beta$ and creates no $B$-dependent additive constant.

The external-coordinate census remains nine for the single finite
square criterion: one input, six point coordinates, one square
parameter $z$, and one encoder coordinate. Its algebra has rank
$64\cdot2=128$ per fixed menu branch; a finite product over menu
branches uses the same external coordinates.

These are soundness criteria. A defining polynomial still requires,
for every nonzero integer input, rational points on the actual source
with its prescribed root label, square parameter, chosen class, and
every retained whole-character guard. That simultaneous realization
is unproved. The height theorem does not establish the required
restricted prime pattern or rational-square incidence. In particular
the established defining-variable bounds remain twenty-five and
thirteen.

\begin{thebibliography}{9}\raggedright
\bibitem{Height}\emph{Signed mixed-power heights and single-cover criteria},
companion manuscript, \nolinkurl{signed_mixed_power_height_v58.tex}.
\bibitem{Conditional}\emph{Thirteen Variables from Direct Elliptic
Index Comparison}, repository manuscript,
\nolinkurl{papers/conditional_few_variables/direct_input13_20261110_v22.tex}.
\bibitem{Single}\emph{A single-cover criterion from a mixed-power
height inequality}, repository manuscript,
\nolinkurl{papers/pole_parameters/single_cover_criterion_20261117_v56.tex}.
\bibitem{Classes}\emph{Variable square classes with controlled height},
repository manuscript,
\nolinkurl{papers/pole_parameters/variable_squareclasses_20261117_v57.tex}.
\end{thebibliography}
\end{document}
